\documentclass[11pt,a4paper]{article}
\usepackage[a4paper,margin=24mm]{geometry}
\usepackage[utf8]{inputenc}
\usepackage[T1]{fontenc}
\usepackage{lmodern}
\usepackage{amsmath,amssymb,amsthm}
\usepackage{booktabs,array,tabularx}
\usepackage{enumitem}
\usepackage{float}
\usepackage{longtable}
\usepackage{microtype}
\usepackage{xcolor}
\definecolor{linkblue}{rgb}{0,0.15,0.5}
\usepackage[colorlinks=true,linkcolor=linkblue,citecolor=linkblue,urlcolor=linkblue]{hyperref}
\usepackage{url}
\allowdisplaybreaks
\emergencystretch=2em

\renewcommand{\thesection}{S\arabic{section}}
\renewcommand{\theequation}{S.\arabic{equation}}
\renewcommand{\thetable}{S\arabic{section}.\arabic{table}}
\theoremstyle{plain}
\newtheorem{theorem}{Theorem}[section]
\newtheorem{proposition}[theorem]{Proposition}
\newtheorem{lemma}[theorem]{Lemma}
\newtheorem{corollary}[theorem]{Corollary}
\theoremstyle{definition}
\newtheorem{definition}[theorem]{Definition}
\newtheorem{example}[theorem]{Example}
\theoremstyle{remark}
\newtheorem{remark}[theorem]{Remark}

\newcommand{\R}{\mathbb{R}}
\newcommand{\Qr}{\mathbb{Q}}
\newcommand{\E}{\mathbb{E}}
\newcommand{\Prob}{\mathbb{P}}
\newcommand{\calA}{\mathcal{A}}
\newcommand{\calK}{\mathcal{K}}
\newcommand{\calM}{\mathcal{M}}
\newcommand{\calP}{\mathcal{P}}
\newcommand{\calR}{\mathcal{R}}
\newcommand{\calS}{\mathcal{S}}
\newcommand{\conv}{\operatorname{conv}}
\newcommand{\aff}{\operatorname{aff}}
\newcommand{\rank}{\operatorname{rank}}
\newcommand{\supp}{\operatorname{supp}}
\newcommand{\capi}[1]{\operatorname{cap}_{#1}}
\newcommand{\ind}[1]{\mathbf{1}[#1]}
\newcommand{\tfr}[2]{\tfrac{#1}{#2}}
\newcommand{\dd}{\,\mathrm{d}}
\newcommand{\mref}[1]{\csname mainlab@#1\endcsname}
\newcommand{\meqref}[1]{(\mref{#1})}
\setlength{\parskip}{2pt}
\renewcommand{\arraystretch}{1.3}

\begin{document}

\title{\textbf{Supplementary material for}\\ \emph{Minimax dual certificates for the common-action obstruction: sparsity, polyhedral recovery, a benchmark law, and the finite review envelope}}
\author{Amin Abaee\\[0.3em]{\small Independent Researcher}\\[0.2em]{\small ORCID: \href{https://orcid.org/0000-0002-0019-1842}{0000-0002-0019-1842}\quad \href{mailto:amin_abaee@ut.ac.ir}{amin\_abaee@ut.ac.ir}}}
\date{October 15, 2026}
\maketitle

\section*{Contents of the supplement}
S1 proof details for the static theory (Sion hypotheses, Helly for infinite families and the rank reduction, the finite convex-hull programme, the polyhedral duality); S2 static instance records (one-dimensional instances, the tight family, the benchmark atoms, margin and demand-floor tables); S3 the two-patch codex enumeration; S4 the random Helly-size study; S5 review-envelope records (the complete two-window tree, the coupling polytope, the sweep with worst couplings, the inherited instances); S6 the integrated certificate and its sign convention; S7 the verification record (check list, Lean scope list, figure provenance); S8 symbol table and sign conventions. Theorem and equation numbers of the main text are quoted as ``main Theorem~\mref{thm:dual}'' etc.; they are read from the compiled main text, and the verification script confirms the correspondence (check TXT.2). The exact module \texttt{p04\_core.py}, the verification script \path{paper04_minimax_dual_certificates_v20_verification.py} with its output, and the figure and table scripts accompany the manuscript.

